\documentclass[11pt,a4paper]{article}

% ── Encodage et langue ──
\usepackage[utf8]{inputenc}
\usepackage[T1]{fontenc}
\usepackage[french]{babel}

% ── Mise en page ──
\usepackage[margin=2.4cm]{geometry}
\usepackage{setspace}
\onehalfspacing

% ── Mathématiques ──
\usepackage{amsmath,amssymb,amsfonts,mathtools}
\usepackage{bm}

% ── Algorithmes ──
\usepackage[ruled,vlined,french,linesnumbered]{algorithm2e}

% ── Figures et couleurs ──
\usepackage{tikz,pgfplots}
\pgfplotsset{compat=1.18}
\usetikzlibrary{arrows.meta,positioning,calc,decorations.pathreplacing,shapes.geometric}
\usepackage{graphicx}
\usepackage{xcolor}

% ── Tableaux ──
\usepackage{booktabs,array,tabularx,colortbl}

% ── Divers ──
\usepackage{enumitem,fancyhdr,titlesec}
\usepackage[colorlinks=true,linkcolor=blue!60!black,citecolor=green!50!black,urlcolor=blue!70!black]{hyperref}

% ── Couleurs du document ──
\definecolor{accent}{HTML}{1A5276}
\definecolor{mygray}{HTML}{F2F3F4}
\definecolor{mygreen}{HTML}{27AE60}
\definecolor{myred}{HTML}{E74C3C}
\definecolor{myorange}{HTML}{E67E22}
\definecolor{myblue}{HTML}{2980B9}

% ── En-têtes ──
\pagestyle{fancy}
\fancyhf{}
\fancyhead[L]{\small\textcolor{accent}{\textsc{RetinAI --- Fondements Mathématiques}}}
\fancyhead[R]{\small\thepage}
\renewcommand{\headrulewidth}{0.4pt}

% ── Titres ──
\titleformat{\section}{\Large\bfseries\color{accent}}{\thesection}{1em}{}[\vspace{-0.5em}\rule{\textwidth}{0.4pt}]
\titleformat{\subsection}{\large\bfseries\color{accent!80!black}}{\thesubsection}{1em}{}

% ── Boîte de définition ──
\usepackage{tcolorbox}
\tcbuselibrary{skins,breakable}
\newtcolorbox{defbox}[1][]{
  colback=mygray, colframe=accent, fonttitle=\bfseries,
  title=#1, breakable, arc=2pt, boxrule=0.8pt
}
\newtcolorbox{notebox}[1][]{
  colback=mygreen!5, colframe=mygreen!70!black, fonttitle=\bfseries,
  title=#1, breakable, arc=2pt, boxrule=0.6pt
}

% ── Commandes personnalisées ──
\newcommand{\R}{\mathbb{R}}
\newcommand{\E}{\mathbb{E}}
\newcommand{\N}{\mathbb{N}}
\DeclareMathOperator*{\argmax}{arg\,max}
\DeclareMathOperator*{\argmin}{arg\,min}
\DeclareMathOperator{\softmax}{softmax}
\DeclareMathOperator{\sigmoid}{sigmoid}
\DeclareMathOperator{\ReLU}{ReLU}

% ══════════════════════════════════════════════════════════════
\begin{document}
\begin{titlepage}
\centering
\vspace*{3cm}
{\Huge\bfseries\color{accent} RetinAI\par}
\vspace{0.8cm}
{\LARGE Fondements Mathématiques et\\Algorithmiques du Pipeline\par}
\vspace{1.5cm}
{\large\textit{Analyse détaillée des algorithmes de Computer Vision\\appliqués à la détection de rétinopathie diabétique}\par}
\vspace{3cm}
{\large Pipeline ML $\cdot$ TensorFlow $\cdot$ Computer Vision\par}
\vspace{1cm}
{\normalsize Mars 2026\par}
\vfill
{\small Document pédagogique --- Toutes les notions sont expliquées de zéro.}
\end{titlepage}

\tableofcontents
\newpage

% ══════════════════════════════════════════════════════════════
\section{Introduction}
% ══════════════════════════════════════════════════════════════

Ce document présente de manière rigoureuse et pédagogique l'ensemble des algorithmes et fondements mathématiques utilisés dans le pipeline \textbf{RetinAI}. Chaque notion est introduite progressivement, des prérequis de base aux formulations avancées, avec des illustrations et des interprétations intuitives.

Le pipeline couvre quatre grandes familles de techniques :

\begin{enumerate}[leftmargin=2em]
  \item \textbf{Prétraitement d'images} : transformations pixel-level pour normaliser et rehausser les images rétiniennes.
  \item \textbf{Réseaux de neurones convolutifs} : architecture profonde pour l'extraction de caractéristiques et la classification.
  \item \textbf{Fonctions de perte} : objectifs d'optimisation adaptés au déséquilibre de classes.
  \item \textbf{Optimisation et inférence} : techniques de compression, augmentation à l'inférence et calibration.
\end{enumerate}

% ══════════════════════════════════════════════════════════════
\section{Prétraitement des images rétiniennes}
% ══════════════════════════════════════════════════════════════

\subsection{Espace colorimétrique et représentation des images}

\begin{defbox}[Image numérique]
Une image couleur est une fonction $I : \{1,\ldots,H\} \times \{1,\ldots,W\} \to \{0,\ldots,255\}^3$ qui associe à chaque pixel $(i,j)$ un triplet $(R,G,B)$ représentant les intensités rouge, verte et bleue. On la représente comme un tenseur $\mathbf{I} \in \R^{H \times W \times 3}$.
\end{defbox}

L'espace \textbf{LAB} (CIE $L^*a^*b^*$) sépare la luminance $L$ (intensité lumineuse) des composantes chromatiques $a$ (axe vert--rouge) et $b$ (axe bleu--jaune). Cette décomposition est exploitée par le CLAHE pour modifier le contraste sans altérer les couleurs.

\subsection{Normalisation Ben-Graham}

\begin{defbox}[Principe]
L'illumination non-uniforme dans les images de fond d'œil introduit un biais spatial. La normalisation Ben-Graham supprime cette composante basse fréquence en soustrayant un flou gaussien local.
\end{defbox}

\subsubsection{Filtre gaussien}

Le noyau gaussien 2D de variance $\sigma^2$ est défini par :
\begin{equation}
  G_\sigma(x,y) = \frac{1}{2\pi\sigma^2}\,\exp\!\left(-\frac{x^2+y^2}{2\sigma^2}\right)
\end{equation}

La convolution de l'image $I$ avec ce noyau produit une version lissée :
\begin{equation}
  I_{\text{blur}}(i,j) = (I * G_\sigma)(i,j) = \sum_{m,n} I(i-m,\,j-n)\,G_\sigma(m,n)
\end{equation}

\subsubsection{Formule de normalisation}

La transformation Ben-Graham s'écrit :
\begin{equation}
  \boxed{I_{\text{norm}}(i,j) = \alpha \cdot I(i,j) - \alpha \cdot I_{\text{blur}}(i,j) + \beta}
\end{equation}
avec typiquement $\alpha = 4$ et $\beta = 128$. Le terme $\alpha(I - I_{\text{blur}})$ isole les hautes fréquences (vaisseaux, micro-anévrismes, exsudats) tandis que $\beta$ recentre les valeurs.

\begin{notebox}[Intuition]
Imaginez que vous regardez une photo avec un éclairage inégal (plus clair à gauche qu'à droite). Le flou gaussien capture cette tendance d'éclairage. En la soustrayant, on obtient une image où seuls les \emph{détails} sont visibles, indépendamment de l'illumination.
\end{notebox}

\subsection{CLAHE --- Contrast Limited Adaptive Histogram Equalization}

\subsubsection{Égalisation d'histogramme classique}

L'histogramme d'une image en niveaux de gris est la distribution des intensités :
\begin{equation}
  h(k) = \#\{(i,j) : I(i,j) = k\}, \quad k \in \{0,\ldots,255\}
\end{equation}

L'égalisation d'histogramme applique la transformation :
\begin{equation}
  T(k) = \left\lfloor 255 \cdot \frac{\text{CDF}(k) - \text{CDF}_{\min}}{H \cdot W - \text{CDF}_{\min}} \right\rfloor
  \quad \text{où} \quad
  \text{CDF}(k) = \sum_{i=0}^{k} h(i)
\end{equation}

\subsubsection{Adaptive Histogram Equalization (AHE)}

L'AHE divise l'image en tuiles (grille $M \times N$, typiquement $8 \times 8$) et applique une égalisation locale à chaque tuile. Pour éviter les artefacts aux frontières, les résultats sont interpolés bilinéairement entre tuiles voisines.

\subsubsection{Limitation du contraste (CLAHE)}

L'AHE peut amplifier le bruit dans les régions homogènes. Le CLAHE résout ce problème en \emph{écrêtant} l'histogramme local :

\begin{algorithm}[H]
\DontPrintSemicolon
\SetAlgoLined
\KwIn{Image $I$, grille $M \times N$, seuil de clip $\tau$}
\KwOut{Image rehaussée $I'$}
Diviser $I$ en tuiles de taille $\frac{H}{M} \times \frac{W}{N}$\;
\Pour{chaque tuile $T_{m,n}$}{
  Calculer l'histogramme local $h_{m,n}$\;
  $\text{excès} \leftarrow \sum_{k} \max(0,\; h_{m,n}(k) - \tau)$\;
  Redistribuer l'excès uniformément : $h'_{m,n}(k) \leftarrow \min(h_{m,n}(k),\,\tau) + \frac{\text{excès}}{256}$\;
  Calculer $\text{CDF}_{m,n}$ à partir de $h'_{m,n}$\;
}
\Pour{chaque pixel $(i,j)$}{
  Identifier les 4 tuiles voisines\;
  $I'(i,j) \leftarrow$ interpolation bilinéaire des 4 CDF appliquées à $I(i,j)$\;
}
\Retour{$I'$}
\caption{CLAHE}
\end{algorithm}

Le paramètre $\tau$ (clip limit) contrôle l'amplification maximale du contraste. Une valeur typique est $\tau = 2.0$.

\subsection{Extraction de la ROI circulaire}

Les images de fond d'œil contiennent un disque rétinien sur fond noir. L'extraction consiste à :

\begin{enumerate}
  \item Seuillage binaire : $B(i,j) = \mathbf{1}[I_{\text{gray}}(i,j) > t]$ avec $t=15$.
  \item Détection du contour le plus grand (disque rétinien).
  \item Calcul du cercle englobant minimal $\mathcal{C}(c_x, c_y, r)$ via l'algorithme de Welzl.
  \item Crop carré centré sur $(c_x, c_y)$ de côté $2(r + m)$ où $m = 0.05 \cdot r$.
  \item Application d'un masque circulaire binaire.
\end{enumerate}

% ══════════════════════════════════════════════════════════════
\section{Réseaux de neurones convolutifs}
% ══════════════════════════════════════════════════════════════

\subsection{Convolution 2D}

\begin{defbox}[Opération de convolution]
Soit une entrée $\mathbf{X} \in \R^{H \times W \times C_{\text{in}}}$ et un filtre $\mathbf{W} \in \R^{k \times k \times C_{\text{in}} \times C_{\text{out}}}$. La convolution produit :
\begin{equation}
  Y_{i,j,f} = \sum_{c=1}^{C_{\text{in}}} \sum_{m=0}^{k-1} \sum_{n=0}^{k-1} W_{m,n,c,f} \cdot X_{i+m,\,j+n,\,c} + b_f
\end{equation}
où $b_f$ est le biais du filtre $f$.
\end{defbox}

\subsection{Fonctions d'activation}

\paragraph{ReLU.} $\ReLU(x) = \max(0, x)$. Introduit la non-linéarité tout en préservant les gradients pour $x > 0$.

\paragraph{Swish / SiLU.} Utilisée dans EfficientNetV2 :
\begin{equation}
  \text{Swish}(x) = x \cdot \sigmoid(\beta x) = \frac{x}{1 + e^{-\beta x}}
\end{equation}
avec $\beta = 1$ par défaut. Plus lisse que ReLU, elle permet un flux de gradient non nul pour $x < 0$.

\paragraph{Softmax.} Pour la couche de sortie (classification) :
\begin{equation}
  \softmax(z_i) = \frac{e^{z_i}}{\sum_{j=1}^{K} e^{z_j}}, \quad i = 1, \ldots, K
\end{equation}
Les sorties forment une distribution de probabilité : $\sum_i \softmax(z_i) = 1$ et $\softmax(z_i) > 0$.

\subsection{Batch Normalization}

Pour chaque mini-batch $\mathcal{B} = \{x_1, \ldots, x_m\}$ :
\begin{align}
  \mu_\mathcal{B} &= \frac{1}{m}\sum_{i=1}^{m} x_i & \sigma_\mathcal{B}^2 &= \frac{1}{m}\sum_{i=1}^{m}(x_i - \mu_\mathcal{B})^2 \\[4pt]
  \hat{x}_i &= \frac{x_i - \mu_\mathcal{B}}{\sqrt{\sigma_\mathcal{B}^2 + \epsilon}} & y_i &= \gamma\,\hat{x}_i + \beta
\end{align}
où $\gamma$ et $\beta$ sont des paramètres appris, et $\epsilon \approx 10^{-5}$ assure la stabilité numérique.

\subsection{EfficientNetV2}

\subsubsection{Scaling composé}

EfficientNet utilise un coefficient composé $\phi$ pour mettre à l'échelle simultanément trois dimensions :
\begin{equation}
  \text{profondeur} : d = \alpha^\phi, \quad
  \text{largeur} : w = \beta^\phi, \quad
  \text{résolution} : r = \gamma^\phi
\end{equation}
sous la contrainte $\alpha \cdot \beta^2 \cdot \gamma^2 \approx 2$ (le FLOPS double pour chaque incrément de $\phi$).

\subsubsection{MBConv et Fused-MBConv}

Le bloc fondamental est le \textbf{Mobile Inverted Bottleneck} (MBConv) :

\begin{enumerate}
  \item \textbf{Expansion} : Conv $1 \times 1$ de $C$ vers $t \cdot C$ canaux (facteur d'expansion $t$).
  \item \textbf{Depthwise convolution} : Conv $k \times k$ par canal (efficace en calcul).
  \item \textbf{Squeeze-and-Excitation} : recalibrage des canaux (voir ci-dessous).
  \item \textbf{Projection} : Conv $1 \times 1$ de $t \cdot C$ vers $C'$ canaux.
  \item \textbf{Connexion résiduelle} : $\mathbf{y} = \mathbf{x} + \text{MBConv}(\mathbf{x})$ si $C = C'$.
\end{enumerate}

EfficientNetV2 remplace les premiers blocs MBConv par des blocs \textbf{Fused-MBConv} qui fusionnent l'expansion et la depthwise convolution en une seule Conv $3 \times 3$ standard, plus rapide sur GPU.

\subsubsection{Squeeze-and-Excitation (SE)}

Le module SE recalibre les canaux en fonction de leur importance globale :

\begin{align}
  \mathbf{s} &= \sigma\!\left(\mathbf{W}_2\,\ReLU\!\left(\mathbf{W}_1\,\text{GAP}(\mathbf{F})\right)\right) \in \R^{C} \\[4pt]
  \widetilde{\mathbf{F}}_c &= s_c \cdot \mathbf{F}_c
\end{align}

où $\text{GAP}(\mathbf{F})_c = \frac{1}{HW}\sum_{i,j} F_{i,j,c}$ est le Global Average Pooling, et $\mathbf{W}_1 \in \R^{C/r \times C}$, $\mathbf{W}_2 \in \R^{C \times C/r}$ avec $r$ le ratio de réduction.

\subsection{Module d'attention spatiale (CBAM)}

Le projet ajoute un module CBAM (Convolutional Block Attention Module) entre le backbone et le classifieur. Il combine attention par canal et attention spatiale.

\subsubsection{Attention par canal}

\begin{equation}
  \mathbf{M}_c(\mathbf{F}) = \sigma\!\Big(\text{MLP}\big(\text{AvgPool}(\mathbf{F})\big) + \text{MLP}\big(\text{MaxPool}(\mathbf{F})\big)\Big)
\end{equation}

Le MLP partagé a une couche cachée de dimension $C/r$ :
\begin{equation}
  \text{MLP}(\mathbf{z}) = \mathbf{W}_2\,\ReLU(\mathbf{W}_1\,\mathbf{z})
\end{equation}

\subsubsection{Attention spatiale}

Après recalibrage des canaux ($\mathbf{F}' = \mathbf{M}_c \odot \mathbf{F}$) :
\begin{equation}
  \mathbf{M}_s(\mathbf{F}') = \sigma\!\Big(f^{7 \times 7}\big([\text{AvgPool}_c(\mathbf{F}');\;\text{MaxPool}_c(\mathbf{F}')]\big)\Big)
\end{equation}

où $\text{AvgPool}_c$ et $\text{MaxPool}_c$ agrègent le long de la dimension des canaux, et $f^{7 \times 7}$ est une convolution $7 \times 7$.

La sortie finale est :
\begin{equation}
  \boxed{\widetilde{\mathbf{F}} = \mathbf{M}_s(\mathbf{F}') \odot \mathbf{F}'}
\end{equation}

\begin{notebox}[Pourquoi l'attention spatiale en imagerie médicale ?]
Les lésions de la rétinopathie diabétique (micro-anévrismes, hémorragies, exsudats) sont des structures \emph{localisées}. Le module d'attention spatiale permet au réseau de se focaliser sur ces régions discriminantes plutôt que de traiter uniformément toute l'image.
\end{notebox}

\subsection{Transfer Learning et Fine-Tuning}

\begin{defbox}[Principe du Transfer Learning]
Un modèle pré-entraîné sur ImageNet (14M images, 21k classes) a appris des représentations visuelles génériques (contours, textures, formes). Ces représentations sont transférables à de nouvelles tâches avec peu de données.
\end{defbox}

\paragraph{Phase 1 --- Warmup.} Le backbone est gelé ($\nabla_{\theta_{\text{backbone}}} = \mathbf{0}$). Seul le classifieur est entraîné avec un learning rate élevé ($\eta = 10^{-3}$).

\paragraph{Phase 2 --- Fine-tuning.} Les derniers blocs du backbone sont dégelés. Un learning rate réduit ($\eta = 10^{-4}$) est utilisé pour ajuster finement les features sans détruire les représentations apprises.

% ══════════════════════════════════════════════════════════════
\section{Fonctions de perte}
% ══════════════════════════════════════════════════════════════

\subsection{Cross-Entropy catégorielle}

Pour $K$ classes, la Cross-Entropy entre la distribution vraie $\mathbf{y}$ (one-hot) et la distribution prédite $\hat{\mathbf{y}}$ est :
\begin{equation}
  \mathcal{L}_{\text{CE}}(\mathbf{y}, \hat{\mathbf{y}}) = -\sum_{k=1}^{K} y_k \log(\hat{y}_k)
\end{equation}

Pour un one-hot $\mathbf{y}$ avec $y_c = 1$ (classe correcte), cela se simplifie en :
\begin{equation}
  \mathcal{L}_{\text{CE}} = -\log(\hat{y}_c) = -\log(p_t)
\end{equation}
où $p_t$ est la probabilité attribuée à la classe correcte.

\subsection{Focal Loss}

\begin{defbox}[Problème du déséquilibre]
Dans le dataset de rétinopathie, le grade 0 (No DR) représente $\sim$73\% des images. La CE classique est dominée par ces exemples faciles et fréquents, au détriment des cas rares mais cliniquement critiques (grades 3--4).
\end{defbox}

La Focal Loss (Lin et al., 2017) module la CE par un facteur $(1 - p_t)^\gamma$ :

\begin{equation}
  \boxed{\mathcal{L}_{\text{FL}}(p_t) = -\alpha_t \cdot (1 - p_t)^\gamma \cdot \log(p_t)}
\end{equation}

\paragraph{Rôle du paramètre $\gamma$ (focusing parameter).}

\begin{center}
\begin{tabular}{lcc}
\toprule
\textbf{Cas} & \textbf{$p_t$} & \textbf{$(1-p_t)^\gamma$ avec $\gamma=2$} \\
\midrule
Bien classé (facile) & 0.95 & $(0.05)^2 = 0.0025$ \\
Modérément classé & 0.70 & $(0.30)^2 = 0.09$ \\
Mal classé (difficile) & 0.20 & $(0.80)^2 = 0.64$ \\
\bottomrule
\end{tabular}
\end{center}

Le facteur focal réduit la contribution des exemples faciles de $\times 400$ par rapport aux exemples difficiles. Le modèle est ainsi forcé de se concentrer sur les cas ambigus.

\paragraph{Pondération par classe $\alpha_t$.}
En plus de la modulation focale, un poids par classe est appliqué :
\begin{equation}
  \alpha_c = \frac{N}{K \cdot N_c}
\end{equation}
où $N$ est le nombre total d'échantillons, $K$ le nombre de classes et $N_c$ le nombre d'échantillons de la classe $c$. Les classes rares reçoivent un poids plus élevé.

\subsection{Label Smoothing}

Au lieu d'utiliser des labels one-hot durs ($y_c \in \{0, 1\}$), le label smoothing remplace :
\begin{equation}
  y'_k = \begin{cases}
    1 - \varepsilon + \frac{\varepsilon}{K} & \text{si } k = c \text{ (classe correcte)} \\[4pt]
    \frac{\varepsilon}{K} & \text{sinon}
  \end{cases}
\end{equation}
avec $\varepsilon = 0.05$. Cela empêche le modèle d'être trop confiant et agit comme un régularisateur.

\subsection{Quadratic Weighted Kappa (QWK) Loss}

Le QWK est la métrique standard pour la classification ordinale de la rétinopathie :

\begin{equation}
  \kappa = 1 - \frac{\sum_{i,j} w_{ij}\,O_{ij}}{\sum_{i,j} w_{ij}\,E_{ij}}
\end{equation}

où :
\begin{itemize}[leftmargin=2em]
  \item $O_{ij}$ : matrice de confusion normalisée (observée).
  \item $E_{ij} = \frac{h_i^{\text{true}} \cdot h_j^{\text{pred}}}{N^2}$ : matrice attendue sous indépendance.
  \item $w_{ij} = \frac{(i-j)^2}{(K-1)^2}$ : matrice de pénalité quadratique.
\end{itemize}

La pénalité quadratique signifie qu'une erreur $0 \to 4$ est pénalisée 16 fois plus qu'une erreur $0 \to 1$, ce qui reflète l'importance clinique.

Pour rendre le QWK différentiable, on utilise une matrice de confusion ``soft'' :
\begin{equation}
  O^{\text{soft}}_{ij} = \frac{1}{N} \sum_{n=1}^{N} y_{n,i} \cdot \hat{y}_{n,j}
\end{equation}
et la loss à minimiser est $\mathcal{L}_{\text{QWK}} = 1 - \kappa$.

% ══════════════════════════════════════════════════════════════
\section{Augmentation des données}
% ══════════════════════════════════════════════════════════════

\subsection{Augmentations géométriques}

Les images de fond d'œil sont \emph{invariantes par rotation} (pas d'orientation anatomique fixe) et par réflexion, ce qui autorise des transformations géométriques agressives :

\begin{itemize}[leftmargin=2em]
  \item \textbf{Rotations} de $0°$, $90°$, $180°$, $270°$ (sans perte d'information).
  \item \textbf{Flip horizontal} et \textbf{vertical}.
\end{itemize}

Cela multiplie par 8 la diversité effective du dataset.

\subsection{MixUp}

MixUp (Zhang et al., 2018) crée des exemples virtuels par combinaison convexe :

\begin{align}
  \tilde{\mathbf{x}} &= \lambda\,\mathbf{x}_i + (1-\lambda)\,\mathbf{x}_j \\
  \tilde{\mathbf{y}} &= \lambda\,\mathbf{y}_i + (1-\lambda)\,\mathbf{y}_j
\end{align}

où $\lambda \sim \text{Beta}(\alpha, \alpha)$ avec $\alpha = 0.2$. La distribution Beta avec un petit $\alpha$ produit des $\lambda$ proches de 0 ou 1, de sorte que les mélanges restent proches d'un des deux exemples originaux.

\begin{notebox}[Effet sur l'espace de décision]
MixUp encourage le modèle à apprendre des frontières de décision \emph{linéaires} entre les classes, ce qui :
\begin{itemize}
  \item Réduit l'overfitting (régularisation implicite).
  \item Améliore la calibration (les probabilités de sortie sont plus fiables).
  \item Renforce la robustesse aux perturbations adversariales.
\end{itemize}
\end{notebox}

\subsection{CutMix}

CutMix (Yun et al., 2019) remplace un patch rectangulaire d'une image par un patch correspondant d'une autre image :

\begin{align}
  \tilde{\mathbf{x}} &= \mathbf{M} \odot \mathbf{x}_i + (\mathbf{1} - \mathbf{M}) \odot \mathbf{x}_j \\
  \tilde{\mathbf{y}} &= \lambda\,\mathbf{y}_i + (1 - \lambda)\,\mathbf{y}_j
\end{align}

où $\mathbf{M} \in \{0,1\}^{H \times W}$ est un masque binaire rectangulaire et $\lambda = 1 - \frac{r_w \cdot r_h}{W \cdot H}$ est la proportion de l'image originale conservée.

La taille du patch est tirée aléatoirement :
\begin{equation}
  r_w = W\sqrt{1-\lambda'}, \quad r_h = H\sqrt{1-\lambda'}, \quad \lambda' \sim \text{Uniform}(0,1)
\end{equation}

% ══════════════════════════════════════════════════════════════
\section{Optimisation de l'entraînement}
% ══════════════════════════════════════════════════════════════

\subsection{Optimiseur AdamW}

AdamW découple la régularisation L2 (weight decay) de l'estimation des moments adaptatifs :

\begin{align}
  \mathbf{m}_t &= \beta_1\,\mathbf{m}_{t-1} + (1 - \beta_1)\,\mathbf{g}_t & \text{(moment d'ordre 1)} \\
  \mathbf{v}_t &= \beta_2\,\mathbf{v}_{t-1} + (1 - \beta_2)\,\mathbf{g}_t^2 & \text{(moment d'ordre 2)} \\
  \hat{\mathbf{m}}_t &= \frac{\mathbf{m}_t}{1 - \beta_1^t}, \quad \hat{\mathbf{v}}_t = \frac{\mathbf{v}_t}{1 - \beta_2^t} & \text{(correction du biais)} \\[6pt]
  \bm{\theta}_t &= \bm{\theta}_{t-1} - \eta\!\left(\frac{\hat{\mathbf{m}}_t}{\sqrt{\hat{\mathbf{v}}_t} + \epsilon} + \lambda_{\text{wd}}\,\bm{\theta}_{t-1}\right) & \text{(mise à jour)}
\end{align}

Le terme $\lambda_{\text{wd}}\,\bm{\theta}_{t-1}$ est le weight decay \emph{découplé} : il pénalise directement la norme des poids, indépendamment de l'adaptation du learning rate.

\subsection{Réduction du learning rate sur plateau}

Quand la loss de validation stagne pendant $p$ epochs (patience), le learning rate est réduit :
\begin{equation}
  \eta_{t+1} = \eta_t \times f, \quad f = 0.5
\end{equation}
avec un plancher $\eta_{\min} = 10^{-7}$.

\subsection{Early Stopping}

L'entraînement est arrêté si la métrique de validation (QWK) ne s'améliore pas pendant $P = 10$ epochs consécutives, et les poids du meilleur epoch sont restaurés. Cela prévient l'overfitting en identifiant automatiquement le point optimal.

% ══════════════════════════════════════════════════════════════
\section{Optimisation du modèle à l'inférence}
% ══════════════════════════════════════════════════════════════

\subsection{Pruning structurel}

Le pruning élimine les poids de faible magnitude :

\begin{equation}
  w'_{ij} = \begin{cases}
    w_{ij} & \text{si } |w_{ij}| > \tau \\
    0 & \text{sinon}
  \end{cases}
\end{equation}

Le seuil $\tau$ est ajusté progressivement via un polynôme pour atteindre le taux de sparsité cible $s_f$ :
\begin{equation}
  s_t = s_f + (s_i - s_f)\left(1 - \frac{t - t_0}{t_1 - t_0}\right)^3
\end{equation}
avec $s_i = 0$ (sparsité initiale) et $s_f = 0.5$ (50\% de poids mis à zéro).

\subsection{Quantification}

La quantification réduit la précision des poids de float32 (32 bits) à int8 (8 bits) :

\begin{equation}
  q = \text{round}\!\left(\frac{w - z}{s}\right), \quad s = \frac{\max(w) - \min(w)}{2^b - 1}
\end{equation}

où $s$ est le facteur d'échelle, $z$ le zero-point et $b = 8$ le nombre de bits. Cela réduit la taille du modèle de $\times 4$ et accélère l'inférence sur CPU.

\subsection{Test-Time Augmentation (TTA)}

Le TTA applique $T$ transformations à chaque image d'entrée et agrège les prédictions :

\begin{equation}
  \boxed{\hat{\mathbf{y}}_{\text{TTA}} = \frac{1}{T} \sum_{t=1}^{T} f_\theta\!\big(\mathcal{A}_t(\mathbf{x})\big)}
\end{equation}

avec $\mathcal{A}_t \in \{\text{id}, R_{90}, R_{180}, R_{270}, H, V, HR_{90}, VR_{90}\}$ (8 transformations).

\paragraph{Estimation de l'incertitude.} L'écart-type entre les prédictions TTA fournit une mesure d'incertitude :
\begin{equation}
  u_k = \sqrt{\frac{1}{T}\sum_{t=1}^{T}\left(\hat{y}_{k}^{(t)} - \bar{y}_k\right)^2}
\end{equation}

Une incertitude élevée signale un cas ambigu nécessitant une attention clinique particulière.

% ══════════════════════════════════════════════════════════════
\section{Évaluation et métriques}
% ══════════════════════════════════════════════════════════════

\subsection{Matrice de confusion}

Pour $K$ classes, la matrice de confusion $\mathbf{C} \in \N^{K \times K}$ est définie par :
\begin{equation}
  C_{ij} = \#\{n : y_n = i \text{ et } \hat{y}_n = j\}
\end{equation}
La diagonale contient les classifications correctes.

\subsection{Métriques par classe}

\begin{align}
  \text{Précision}_k &= \frac{C_{kk}}{\sum_{i} C_{ik}} & \text{(parmi les prédits $k$, combien sont vrais $k$ ?)} \\[6pt]
  \text{Rappel}_k &= \frac{C_{kk}}{\sum_{j} C_{kj}} & \text{(parmi les vrais $k$, combien sont détectés ?)} \\[6pt]
  F_1^{(k)} &= 2 \cdot \frac{\text{Précision}_k \cdot \text{Rappel}_k}{\text{Précision}_k + \text{Rappel}_k}
\end{align}

\subsection{AUC-ROC (Area Under the ROC Curve)}

Pour chaque classe $k$ (one-vs-rest), la courbe ROC trace le taux de vrais positifs en fonction du taux de faux positifs en faisant varier le seuil de décision $\tau$ :

\begin{align}
  \text{TPR}(\tau) &= \frac{\text{TP}(\tau)}{\text{TP}(\tau) + \text{FN}(\tau)} \\[4pt]
  \text{FPR}(\tau) &= \frac{\text{FP}(\tau)}{\text{FP}(\tau) + \text{TN}(\tau)}
\end{align}

L'AUC est l'aire sous cette courbe :
\begin{equation}
  \text{AUC}_k = \int_0^1 \text{TPR}_k\big(\text{FPR}_k^{-1}(x)\big)\,dx
\end{equation}

Un AUC de 1.0 indique un classifieur parfait, 0.5 un classifieur aléatoire. L'AUC macro est la moyenne non pondérée des AUC par classe.

\subsection{Expected Calibration Error (ECE)}

L'ECE mesure l'écart entre la confiance du modèle et sa précision réelle :

\begin{equation}
  \text{ECE} = \sum_{b=1}^{B} \frac{|S_b|}{N} \left|\text{acc}(S_b) - \text{conf}(S_b)\right|
\end{equation}

où $S_b$ est l'ensemble des échantillons dont la confiance tombe dans le bin $b$, et :
\begin{equation}
  \text{acc}(S_b) = \frac{1}{|S_b|}\sum_{i \in S_b} \mathbf{1}[\hat{y}_i = y_i], \quad
  \text{conf}(S_b) = \frac{1}{|S_b|}\sum_{i \in S_b} \max_k \hat{y}_{i,k}
\end{equation}

Un modèle parfaitement calibré a $\text{ECE} = 0$.

% ══════════════════════════════════════════════════════════════
\section{Explicabilité : Grad-CAM}
% ══════════════════════════════════════════════════════════════

\begin{defbox}[Gradient-weighted Class Activation Mapping]
Grad-CAM produit une carte de chaleur montrant quelles régions spatiales de l'image contribuent le plus à la décision du modèle pour une classe donnée.
\end{defbox}

\subsection{Formulation}

Soit $A^k_{ij}$ les feature maps de la dernière couche convolutive (canal $k$, position spatiale $(i,j)$) et $y^c$ le score de la classe cible $c$ \emph{avant} softmax.

\paragraph{Étape 1.} Calcul des poids d'importance par canal via les gradients :
\begin{equation}
  \alpha_k^c = \underbrace{\frac{1}{Z}\sum_{i}\sum_{j}}_{\text{GAP}} \frac{\partial y^c}{\partial A^k_{ij}}
\end{equation}

\paragraph{Étape 2.} Combinaison linéaire pondérée suivie d'une ReLU :
\begin{equation}
  \boxed{L^c_{\text{Grad-CAM}}(i,j) = \ReLU\!\left(\sum_{k} \alpha_k^c \cdot A^k_{ij}\right)}
\end{equation}

La ReLU ne conserve que les contributions \emph{positives} (régions qui augmentent le score de la classe $c$).

\paragraph{Étape 3.} Upsampling bilinéaire à la résolution de l'image d'entrée et normalisation dans $[0, 1]$.

\begin{notebox}[Importance clinique]
En imagerie médicale, l'explicabilité n'est pas un luxe mais une nécessité. Un radiologue ne peut pas faire confiance à un modèle ``boîte noire''. Les cartes Grad-CAM permettent de vérifier que le modèle regarde les bonnes zones (micro-anévrismes, hémorragies, néovascularisation) et non des artefacts (bords de l'image, reflets).
\end{notebox}

% ══════════════════════════════════════════════════════════════
\section{Récapitulatif des notations}
% ══════════════════════════════════════════════════════════════

\begin{center}
\renewcommand{\arraystretch}{1.3}
\begin{tabularx}{\textwidth}{lX}
\toprule
\textbf{Symbole} & \textbf{Signification} \\
\midrule
$\mathbf{I} \in \R^{H \times W \times 3}$ & Image d'entrée (hauteur $H$, largeur $W$, 3 canaux RGB) \\
$G_\sigma$ & Noyau gaussien de variance $\sigma^2$ \\
$K$ & Nombre de classes ($K = 5$ : grades 0 à 4) \\
$p_t$ & Probabilité assignée à la classe correcte \\
$\gamma$ & Paramètre de focalisation de la Focal Loss \\
$\alpha_t$ & Poids de la classe $t$ dans la Focal Loss \\
$\kappa$ & Quadratic Weighted Kappa \\
$\lambda$ & Coefficient de mélange MixUp / CutMix \\
$\eta$ & Learning rate \\
$\bm{\theta}$ & Paramètres du modèle \\
$\mathbf{m}_t, \mathbf{v}_t$ & Moments d'ordre 1 et 2 (AdamW) \\
$s, z$ & Facteur d'échelle et zero-point (quantification) \\
$A^k_{ij}$ & Feature map du canal $k$ à la position $(i,j)$ \\
$\alpha_k^c$ & Poids Grad-CAM du canal $k$ pour la classe $c$ \\
\bottomrule
\end{tabularx}
\end{center}

\vfill
\begin{center}
\textcolor{accent}{\rule{0.5\textwidth}{0.5pt}}\\[0.5em]
\textit{Fin du document --- Fondements mathématiques et algorithmiques}
\end{center}

\end{document}
